\begin{SideNotePage}{
  \textbf{Top (Voting methods):} Ranked-choice voting is a system in which voters express their preferences by submitting complete rankings of all candidates, and the system aggregates them into a ranked list or a single winner. Different aggregation methods (Borda count, instant-runoff voting (IRV), plurality, Condorcet) can produce distinct winners from identical voter rankings, demonstrating the inherent ambiguity in collective decision-making. The same preference profile can yield different outcomes depending on which features of the rankings are emphasised by the chosen method. This indeterminacy means there is no canonical way to translate individual preferences into collective choices. \par

  \textbf{Bottom (Arrow's Theorem):} Arrow's impossibility theorem proves that no ranked-choice voting method can satisfy all four fairness criteria simultaneously when there are at least three alternatives and two voters. Every democratic aggregation method must compromise at least one criterion, making trade-offs unavoidable in social choice. \par
}{09_ArrowTheoremTopology/09_ Real Democracy Has Never Been Tried.pdf}
\end{SideNotePage}
